% from ega0-1-fr.tex
\oldpage[0\textsubscript{I}]{40}
qui ne sera en général ni injectif ni surjectif. Si $\mathcal{P}$ est un
troisième $\mathcal{A}$-Module, on vérifie aussitôt que le diagramme
\[
\begin{tikzcd}[column sep=large,row sep=large]
\Psi_*(\mathcal{M})\otimes_{\mathcal{B}}\Psi_*(\mathcal{N})
  \otimes_{\mathcal{B}}\Psi_*(\mathcal{P})
  \arrow[r]
  \arrow[d]
&
\Psi_*(\mathcal{M}\otimes_{\mathcal{A}}\mathcal{N})
  \otimes_{\mathcal{B}}\Psi_*(\mathcal{P})
  \arrow[d]
\\
\Psi_*(\mathcal{M})\otimes_{\mathcal{B}}
  \Psi_*(\mathcal{N}\otimes_{\mathcal{A}}\mathcal{P})
  \arrow[r]
&
\Psi_*(\mathcal{M}\otimes_{\mathcal{A}}\mathcal{N}
  \otimes_{\mathcal{A}}\mathcal{P})
\end{tikzcd}
\tag{4.2.2.2}\label{0.4.2.2.2-fr}
\]
est commutatif.
\end{env}

\begin{env}[4.2.3]
\label{0.4.2.3-fr}
Soient $\mathcal{M}$, $\mathcal{N}$ deux $\mathcal{A}$-Modules. Pour tout
ouvert $U\subset Y$, on a par définition
\[
  \Gamma\bigl(\psi^{-1}(U),
    \mathcal{H}\!om_{\mathcal{A}}(\mathcal{M},\mathcal{N})\bigr)
  =
  \operatorname{Hom}_{\mathcal{A}|V}
    (\mathcal{M}|V,\mathcal{N}|V),
\]
où on a posé pour alléger $V=\psi^{-1}(U)$~; l'application
$u\to\Psi_*(u)$ est un homomorphisme
\[
  \operatorname{Hom}_{\mathcal{A}|V}(\mathcal{M}|V,\mathcal{N}|V)
  \longrightarrow
  \operatorname{Hom}_{\mathcal{B}|U}
    (\Psi_*(\mathcal{M})|U,\Psi_*(\mathcal{N})|U)
\]
pour les structures de $\Gamma(U,\mathcal{B})$-modules~; ces homomorphismes
étant compatibles aux opérations de restriction, définissent donc un
homomorphisme canonique fonctoriel de $\mathcal{B}$-Modules
\[
  \Psi_*\bigl(\mathcal{H}\!om_{\mathcal{A}}(\mathcal{M},\mathcal{N})\bigr)
  \longrightarrow
  \mathcal{H}\!om_{\mathcal{B}}
    (\Psi_*(\mathcal{M}),\Psi_*(\mathcal{N})).
  \tag{4.2.3.1}\label{0.4.2.3.1-fr}
\]
\end{env}

\begin{env}[4.2.4]
\label{0.4.2.4-fr}
Si $\mathcal{C}$ est une $\mathcal{A}$-Algèbre, l'homomorphisme composé
\[
  \Psi_*(\mathcal{C})\otimes_{\mathcal{B}}\Psi_*(\mathcal{C})
  \longrightarrow
  \Psi_*(\mathcal{C}\otimes_{\mathcal{A}}\mathcal{C})
  \longrightarrow
  \Psi_*(\mathcal{C})
\]
définit sur $\Psi_*(\mathcal{C})$ une structure de $\mathcal{B}$-Algèbre,
comme il résulte de (4.2.2.2). On voit de même que si $\mathcal{M}$ est un
$\mathcal{C}$-Module, $\Psi_*(\mathcal{M})$ est muni canoniquement d'une
structure de $\Psi_*(\mathcal{C})$-Module.
\end{env}

\begin{env}[4.2.5]
\label{0.4.2.5-fr}
Considérons en particulier le cas où $X$ est un sous-espace \emph{fermé} de
$Y$ et où $\psi$ est l'injection canonique $j:X\to Y$. Si
$\mathcal{B}'=\mathcal{B}|X=j^*(\mathcal{B})$ est la restriction du faisceau
d'anneaux $\mathcal{B}$ à $X$, un $\mathcal{A}$-Module $\mathcal{M}$ peut
être considéré comme $\mathcal{B}'$-Module au moyen de l'homomorphisme
$\theta^\sharp:\mathcal{B}'\to\mathcal{A}$~; $\Psi_*(\mathcal{M})$ est alors
le $\mathcal{B}$-Module qui induit $\mathcal{M}$ sur $X$ et $0$ ailleurs. Si
$\mathcal{N}$ est un second $\mathcal{A}$-Module,
$\Psi_*(\mathcal{M})\otimes_{\mathcal{B}}\Psi_*(\mathcal{N})$ s'identifie
donc à $\Psi_*(\mathcal{M}\otimes_{\mathcal{B}'}\mathcal{N})$ et
$\mathcal{H}\!om_{\mathcal{B}}(\Psi_*(\mathcal{M}),\Psi_*(\mathcal{N}))$
à
$\Psi_*\bigl(\mathcal{H}\!om_{\mathcal{B}'}(\mathcal{M},\mathcal{N})\bigr)$.
\end{env}

\begin{env}[4.2.6]
\label{0.4.2.6-fr}
Soient $(Z,\mathcal{C})$ un troisième espace annelé,
$\Psi'=(\psi',\theta')$ un morphisme
$(Y,\mathcal{B})\to(Z,\mathcal{C})$~; si $\Psi''$ est le morphisme composé
$\Psi'\circ\Psi$, il est clair que l'on a
$\Psi''_*=\Psi'_*\circ\Psi_*$.
\end{env}

